% !TEX TS-program = lualatex
% !TEX encoding = UTF-8 Unicode
\documentclass[11pt,a4paper]{article}

\usepackage{fontspec}
\usepackage[italian]{babel}
\usepackage{geometry}
\usepackage{enumitem}
\usepackage{xcolor}
\usepackage{hyperref}
\usepackage{listings}
\usepackage{titlesec}

\geometry{
	margin=2.2cm
}

\setmainfont{Latin Modern Roman}
\setsansfont{Latin Modern Sans}
\setmonofont{Latin Modern Mono}

\hypersetup{
	colorlinks=true,
	linkcolor=blue,
	urlcolor=blue
}

\setlist[itemize]{itemsep=2pt, topsep=4pt}
\setlist[enumerate]{itemsep=2pt, topsep=4pt}

\lstdefinestyle{codice}{
	basicstyle=\ttfamily\small,
	breaklines=true,
	frame=single,
	columns=fullflexible,
	backgroundcolor=\color{gray!8}
}

\title{Istruzioni per il programma ESP32 FabGL WiFi RSS}
\author{}
\date{}

\begin{document}
	
	\maketitle
	
	\section{Descrizione generale}
	
	Questo programma permette di utilizzare un ESP32 collegato a un'uscita VGA tramite la libreria FabGL per visualizzare notizie provenienti da feed RSS.
	
	Il dispositivo si collega alla rete WiFi, sincronizza l'orologio tramite NTP e mostra sullo schermo:
	
	\begin{itemize}
		\item la data e l'ora locale;
		\item una barra di avanzamento che indica il tempo mancante al prossimo aggiornamento RSS;
		\item il nome del feed RSS attualmente selezionato;
		\item l'elenco delle notizie più recenti;
		\item la data e l'ora delle singole notizie;
		\item i titoli delle notizie con ritorno a capo automatico.
	\end{itemize}
	
	Il programma non utilizza più un file \texttt{config.h}: le credenziali WiFi e l'elenco dei feed RSS possono essere configurati tramite una pagina web integrata.
	
	\section{Librerie richieste}
	
	Per compilare correttamente il programma sono necessarie le seguenti librerie:
	
	\begin{itemize}
		\item \texttt{FabGL}, versione 1.0.9;
		\item \texttt{Bounce2};
		\item librerie standard ESP32:
		\begin{itemize}
			\item \texttt{WiFi.h};
			\item \texttt{WiFiClientSecure.h};
			\item \texttt{HTTPClient.h};
			\item \texttt{time.h}.
		\end{itemize}
		\item il modulo locale \texttt{LkWiFiConfigPortal.h}.
	\end{itemize}
	
	Il file \texttt{LkWiFiConfigPortal.h} deve essere presente nella cartella del progetto insieme al file principale \texttt{.ino}.
	
	\section{Impostazioni di compilazione}
	
	Usare Arduino IDE con il core ESP32 di Espressif.
	
	\subsection{Versioni consigliate}
	
	\begin{itemize}
		\item Core ESP32 by Espressif Systems: \textbf{2.0.4};
		\item Scheda: \textbf{ESP32 Dev Module};
		\item FabGL: \textbf{1.0.9}.
	\end{itemize}
	
	\subsection{Parametri della scheda}
	
	Impostare Arduino IDE nel seguente modo:
	
	\begin{lstlisting}[style=codice]
		Board: ESP32 Dev Module
		Upload Speed: 921600 oppure 115200
		CPU Frequency: 240 MHz
		Flash Frequency: 80 MHz
		Flash Mode: QIO
		Flash Size: 4MB
		Partition Scheme: Default 4MB with spiffs
		PSRAM: Disabled
		Core Debug Level: None
		Port: la porta USB del tuo ESP32
	\end{lstlisting}
	
	\section{Collegamento del pulsante RSS}
	
	Il programma usa un pulsante collegato al pin:
	
	\begin{lstlisting}[style=codice]
		GPIO 33
	\end{lstlisting}
	
	Nel codice il pin è definito così:
	
	\begin{lstlisting}[style=codice]
		const int BTN_NEXT_RSS = 33;
	\end{lstlisting}
	
	Il pulsante deve essere collegato in modo che, quando viene premuto, porti il pin a massa.
	
	Il programma usa la resistenza interna di pull-up dell'ESP32:
	
	\begin{lstlisting}[style=codice]
		pinMode(BTN_NEXT_RSS, INPUT_PULLUP);
	\end{lstlisting}
	
	Quindi lo schema logico è:
	
	\begin{itemize}
		\item pulsante rilasciato: ingresso alto;
		\item pulsante premuto: ingresso basso.
	\end{itemize}
	
	\section{Modalità di configurazione}
	
	Il programma può partire in due modalità:
	
	\begin{itemize}
		\item modalità normale;
		\item modalità configurazione.
	\end{itemize}
	
	\subsection{Avvio normale}
	
	Se il pulsante RSS non è premuto all'accensione, il dispositivo tenta di collegarsi alla rete WiFi già configurata.
	
	Se la connessione WiFi è disponibile, il programma:
	
	\begin{enumerate}
		\item sincronizza l'orologio tramite NTP;
		\item scarica il feed RSS selezionato;
		\item mostra data, ora e notizie sullo schermo VGA.
	\end{enumerate}
	
	\subsection{Avvio in modalità configurazione}
	
	Per avviare la modalità configurazione:
	
	\begin{enumerate}
		\item spegnere l'ESP32;
		\item tenere premuto il pulsante RSS;
		\item accendere l'ESP32;
		\item rilasciare il pulsante dopo l'avvio.
	\end{enumerate}
	
	In questa modalità il dispositivo crea un Access Point WiFi.
	
	Sullo schermo compare un messaggio simile a:
	
	\begin{lstlisting}[style=codice]
		Modalita configurazione: AP ESP32-RSS-SETUP,
		pagina http://192.168.4.1/config.html
	\end{lstlisting}
	
	Collegarsi quindi alla rete WiFi creata dall'ESP32 e aprire nel browser:
	
	\begin{center}
		\url{http://192.168.4.1/config.html}
	\end{center}
	
	Da questa pagina è possibile configurare le credenziali WiFi.
	
	\section{Configurazione dei feed RSS}
	
	Il programma contiene alcuni feed RSS predefiniti:
	
	\begin{itemize}
		\item Il Fatto Quotidiano;
		\item The Guardian;
		\item New York Times;
		\item ANSA Cronaca;
		\item ANSA Home;
		\item ANSA Mondo.
	\end{itemize}
	
	Questi feed sono usati come valori iniziali, ma possono essere modificati tramite pagina web.
	
	La pagina per modificare l'elenco dei feed RSS è:
	
	\begin{center}
		\url{http://IP_ESP32/rss.html}
	\end{center}
	
	Dove \texttt{IP\_ESP32} è l'indirizzo IP assegnato dal router all'ESP32.
	
	Esempio:
	
	\begin{center}
		\url{http://192.168.1.50/rss.html}
	\end{center}
	
	\section{Cambio del feed RSS}
	
	Il pulsante collegato al pin GPIO 33 serve anche per cambiare il feed RSS visualizzato.
	
	A ogni pressione del pulsante, il programma passa al feed successivo.
	
	Quando arriva all'ultimo feed configurato, torna automaticamente al primo.
	
	Per evitare cambi accidentali o rimbalzi del pulsante, il programma usa:
	
	\begin{itemize}
		\item la libreria \texttt{Bounce2};
		\item un intervallo di debounce di 80 ms;
		\item un tempo di stabilizzazione della selezione RSS di 2 secondi.
	\end{itemize}
	
	Il tempo di stabilizzazione è definito da:
	
	\begin{lstlisting}[style=codice]
		const unsigned long RSS_SELECTION_SETTLE_MS = 2000;
	\end{lstlisting}
	
	Questo significa che, dopo il cambio del feed, il programma aspetta 2 secondi prima di scaricare effettivamente le nuove notizie.
	
	\section{Aggiornamento automatico delle notizie}
	
	Le notizie vengono aggiornate automaticamente ogni minuto.
	
	Il tempo di aggiornamento è definito da:
	
	\begin{lstlisting}[style=codice]
		const uint32_t FETCH_REFRESH_MS = 1UL * 60UL * 1000UL;
	\end{lstlisting}
	
	Per modificare l'intervallo di aggiornamento, cambiare questo valore.
	
	Per esempio, per aggiornare ogni 5 minuti:
	
	\begin{lstlisting}[style=codice]
		const uint32_t FETCH_REFRESH_MS = 5UL * 60UL * 1000UL;
	\end{lstlisting}
	
	\section{Barra di avanzamento}
	
	Sulla prima riga dello schermo viene visualizzata una piccola barra di avanzamento che indica quanto tempo è passato dall'ultimo aggiornamento RSS.
	
	La lunghezza della barra è definita da:
	
	\begin{lstlisting}[style=codice]
		const int RSS_BAR_LEN = 20;
	\end{lstlisting}
	
	La barra si riempie progressivamente fino al successivo aggiornamento automatico delle notizie.
	
	\section{Numero massimo di notizie}
	
	Il programma gestisce al massimo 20 notizie per ogni feed RSS.
	
	Il valore è definito da:
	
	\begin{lstlisting}[style=codice]
		const int MAX_NEWS = 20;
	\end{lstlisting}
	
	Se il feed contiene più di 20 notizie, vengono considerate solo le prime 20 trovate e ordinate in base alla data.
	
	\section{Sincronizzazione dell'orologio}
	
	Quando il WiFi è disponibile, il programma sincronizza l'orario tramite NTP usando i server:
	
	\begin{itemize}
		\item \texttt{pool.ntp.org};
		\item \texttt{time.google.com}.
	\end{itemize}
	
	Il fuso orario è impostato per l'Italia:
	
	\begin{lstlisting}[style=codice]
		const char* TZ_EUROPE_ROME =
		"CET-1CEST,M3.5.0/2,M10.5.0/3";
	\end{lstlisting}
	
	Questo permette di gestire automaticamente ora solare e ora legale.
	
	\section{Visualizzazione sullo schermo}
	
	La riga 0 dello schermo mostra:
	
	\begin{itemize}
		\item giorno della settimana;
		\item giorno del mese;
		\item mese;
		\item anno;
		\item ora;
		\item barra di avanzamento RSS;
		\item nome del feed selezionato.
	\end{itemize}
	
	Le righe successive mostrano le notizie.
	
	Per ogni notizia vengono visualizzati:
	
	\begin{itemize}
		\item la data e l'ora della notizia;
		\item il titolo della notizia.
	\end{itemize}
	
	I titoli troppo lunghi vengono spezzati automaticamente su più righe.
	
	\section{Colori delle notizie}
	
	Ogni titolo viene visualizzato con una combinazione di colori calcolata in base al testo della notizia.
	
	Il programma evita alcune combinazioni poco leggibili, per esempio:
	
	\begin{itemize}
		\item giallo su ciano;
		\item ciano su giallo;
		\item verde su ciano;
		\item rosso su magenta;
		\item bianco su giallo;
		\item bianco su verde;
		\item bianco su ciano.
	\end{itemize}
	
	Questo permette di differenziare visivamente le notizie mantenendo una buona leggibilità.
	
	\section{Pulizia dei titoli RSS}
	
	Prima di mostrare i titoli, il programma esegue alcune operazioni di pulizia:
	
	\begin{itemize}
		\item rimuove eventuali sezioni CDATA;
		\item converte alcune entità HTML;
		\item normalizza apostrofi, virgolette e trattini;
		\item converte alcuni caratteri accentati per la visualizzazione corretta sul terminale FabGL;
		\item rimuove parole come \texttt{DIRETTA}, \texttt{LIVE} e \texttt{VIDEO}.
	\end{itemize}
	
	Questa pulizia rende i titoli più adatti alla visualizzazione su schermo VGA testuale.
	
	\section{Gestione della connessione WiFi}
	
	Durante il ciclo principale, il programma controlla continuamente lo stato del portale WiFi.
	
	Se il WiFi viene configurato dopo l'avvio, il programma:
	
	\begin{enumerate}
		\item sincronizza l'orologio;
		\item forza un primo download delle notizie;
		\item aggiorna la schermata.
	\end{enumerate}
	
	Questo consente di configurare il dispositivo senza dover necessariamente ricompilare il firmware.
	
	\section{Funzionamento del ciclo principale}
	
	Nel \texttt{loop()} il programma esegue continuamente queste operazioni:
	
	\begin{enumerate}
		\item gestisce le richieste web del portale di configurazione;
		\item verifica se il WiFi è stato appena configurato;
		\item controlla se la configurazione RSS è stata modificata da web;
		\item legge il pulsante per il cambio feed;
		\item applica l'eventuale cambio feed dopo il tempo di stabilizzazione;
		\item aggiorna la riga dell'orologio ogni secondo;
		\item aggiorna periodicamente le notizie RSS.
	\end{enumerate}
	
	L'orologio viene aggiornato senza ridisegnare tutta la schermata, così si evita lo sfarfallamento.
	
	\section{Procedura consigliata di utilizzo}
	
	\begin{enumerate}
		\item Caricare il programma sull'ESP32.
		\item Collegare il sistema VGA secondo lo schema previsto da FabGL.
		\item Collegare il pulsante RSS al GPIO 33 e a massa.
		\item Accendere l'ESP32 tenendo premuto il pulsante per entrare in configurazione.
		\item Collegarsi all'Access Point creato dall'ESP32.
		\item Aprire:
		\begin{center}
			\url{http://192.168.4.1/config.html}
		\end{center}
		\item Inserire le credenziali WiFi.
		\item Riavviare l'ESP32 in modalità normale.
		\item Quando l'ESP32 è collegato al router, aprire:
		\begin{center}
			\url{http://IP_ESP32/rss.html}
		\end{center}
		\item Modificare eventualmente l'elenco dei feed RSS.
		\item Usare il pulsante per cambiare feed RSS.
	\end{enumerate}
	
	\section{Personalizzazioni principali}
	
	\subsection{Cambiare il tempo di aggiornamento RSS}
	
	Modificare:
	
	\begin{lstlisting}[style=codice]
		const uint32_t FETCH_REFRESH_MS = 1UL * 60UL * 1000UL;
	\end{lstlisting}
	
	\subsection{Cambiare il numero massimo di notizie}
	
	Modificare:
	
	\begin{lstlisting}[style=codice]
		const int MAX_NEWS = 20;
	\end{lstlisting}
	
	\subsection{Cambiare il pin del pulsante}
	
	Modificare:
	
	\begin{lstlisting}[style=codice]
		const int BTN_NEXT_RSS = 33;
	\end{lstlisting}
	
	\subsection{Cambiare i feed RSS predefiniti}
	
	Modificare l'array:
	
	\begin{lstlisting}[style=codice]
		const LkRSSFeed DEFAULT_RSS_FEEDS[] = {
			{ "IL FATTO QUOTIDIANO ", "https://www.ilfattoquotidiano.it/feed/"},
			{ "THE GUARDIAN", "https://www.theguardian.com/world/rss"},
			{ "NEW YORK TIMES", "https://rss.nytimes.com/services/xml/rss/nyt/World.xml"},
			{ "ANSA CRONACA", "https://www.ansa.it/sito/notizie/cronaca/cronaca_rss.xml" },
			{ "ANSA HOME",    "https://www.ansa.it/sito/ansait_rss.xml" },
			{ "ANSA MONDO",   "https://www.ansa.it/sito/notizie/mondo/mondo_rss.xml" },
		};
	\end{lstlisting}
	
	Questi valori sono solo i valori iniziali: dopo la configurazione tramite pagina web, il programma userà i feed salvati nel portale.
	
	\section{Note importanti}
	
	\begin{itemize}
		\item Il programma usa connessioni HTTPS ma disabilita la verifica del certificato tramite \texttt{client.setInsecure()}.
		\item Questo semplifica il download dei feed RSS da ESP32, ma non offre una verifica crittografica completa del certificato.
		\item I feed RSS devono essere raggiungibili dall'ESP32 tramite la rete WiFi configurata.
		\item Alcuni feed potrebbero non funzionare se il server richiede compressione, redirect particolari o formati XML non compatibili.
		\item Il programma forza l'header \texttt{Accept-Encoding: identity} per evitare risposte compresse difficili da gestire sull'ESP32.
	\end{itemize}
	
	\section{Risoluzione dei problemi}
	
	\subsection{Non compare nessuna notizia}
	
	Controllare che:
	
	\begin{itemize}
		\item l'ESP32 sia connesso al WiFi;
		\item l'indirizzo del feed RSS sia corretto;
		\item il feed RSS sia raggiungibile da browser;
		\item il router permetta l'accesso a Internet;
		\item il feed restituisca codice HTTP 200.
	\end{itemize}
	
	\subsection{L'ora non viene visualizzata}
	
	Controllare che:
	
	\begin{itemize}
		\item il WiFi sia connesso;
		\item l'ESP32 possa raggiungere i server NTP;
		\item la rete non blocchi il traffico NTP.
	\end{itemize}
	
	\subsection{Il pulsante non cambia feed}
	
	Controllare che:
	
	\begin{itemize}
		\item il pulsante sia collegato tra GPIO 33 e GND;
		\item il pin configurato nel codice sia corretto;
		\item il pulsante funzioni elettricamente;
		\item siano presenti almeno due feed RSS configurati.
	\end{itemize}
	
	\subsection{Caratteri accentati visualizzati male}
	
	Il programma converte alcuni caratteri accentati italiani in byte compatibili con la visualizzazione terminale.
	
	Se qualche carattere viene visualizzato ancora in modo errato, occorre aggiungere una nuova conversione nella funzione:
	
	\begin{lstlisting}[style=codice]
		toDisplayCharset()
	\end{lstlisting}
	
	\section{Conclusione}
	
	Il programma realizza un visualizzatore RSS autonomo basato su ESP32 e FabGL.
	
	La configurazione tramite pagina web permette di modificare WiFi e feed RSS senza ricompilare il firmware. Il pulsante fisico consente di cambiare rapidamente sorgente RSS, mentre la schermata VGA mostra data, ora e notizie aggiornate automaticamente.
\end{document}
```
